\section{Z-Function and Pattern Matching}
\label{sec:str-z-function}

\problemstatement{Compute the Z-function of a string -- for each index, the
length of the longest substring starting there that matches a prefix of the
whole string -- and use it to find all occurrences of a pattern in a text.}

\begin{patternbox}
The \textbf{Z-function} measures self-similarity in linear time using a
maintained ``Z-box'' (the rightmost prefix-match interval). For pattern
matching, run it on $\textit{pattern}+\#+\textit{text}$: positions in the
text part with $Z=|\textit{pattern}|$ are matches.
\end{patternbox}

\subsection*{The Z-box makes it linear}

$Z[i]$ is how far the substring at $i$ agrees with the prefix. Maintain the
interval $[l,r]$ of the rightmost prefix match found so far. If $i$ lies
inside $[l,r]$, the mirror value $Z[i-l]$ gives a free lower bound for
$Z[i]$ (capped by $r-i$); then extend by direct comparison, moving $r$
forward. Each character is compared a constant number of times amortised, so
the whole array is $O(n)$.

\begin{ideabox}[Reuse the Rightmost Match to Skip Comparisons]
Inside the current Z-box, a position mirrors an earlier one whose Z-value is
known, so most of $Z[i]$ is free. Only growth past the box costs
comparisons, and those advance the box -- bounding total work linearly.
\end{ideabox}

\subsection*{Algorithm}

\begin{enumerate}
  \item Form $w=\textit{pattern}+\#+\textit{text}$; compute its Z-array
        with the Z-box recurrence.
  \item Every index $i$ in the text region with $Z[i]=|\textit{pattern}|$
        is an occurrence.
\end{enumerate}

\subsection*{Dry run}

\dryrun{pattern \texttt{"ab"}, text \texttt{"abcab"}; $w=$\texttt{"ab\#abcab"}.}

\begin{itemize}
  \item $Z=2$ at the two positions where \texttt{"ab"} begins in the text.
  \item Matches at text indices $0$ and $3$.
\end{itemize}

\subsection*{C++ implementation}

\begin{lstlisting}
#include <bits/stdc++.h>
using namespace std;

vector<int> zFunction(const string& w) {
    int n = w.size();
    vector<int> z(n, 0);
    z[0] = n;
    int l = 0, r = 0;
    for (int i = 1; i < n; i++) {
        if (i < r) z[i] = min(r - i, z[i - l]);    // mirror inside the Z-box
        while (i + z[i] < n && w[z[i]] == w[i + z[i]]) z[i]++;
        if (i + z[i] > r) { l = i; r = i + z[i]; }
    }
    return z;
}

int main() {
    string pat = "ab", text = "abcab";
    string w = pat + "#" + text;
    vector<int> z = zFunction(w);
    int m = pat.size();
    for (int i = m + 1; i < (int)w.size(); i++)
        if (z[i] == m) cout << i - m - 1 << " ";   // text index
    cout << "\n";                                  // 0 3
    return 0;
}
\end{lstlisting}

\begin{complexitybox}
\textbf{Time Complexity.} $O(n+m)$. \textbf{Space Complexity.} $O(n+m)$.
\end{complexitybox}

\begin{takeawaybox}
The Z-function computes prefix self-matches in linear time via the Z-box,
and pattern matching drops out by prepending the pattern with a separator.
It is interchangeable with KMP; keep both as linear matchers.
\end{takeawaybox}
